% Principal GDZ PPN236018647, native464–465; HKUST465–466.
\clearpage\SourceLeaf{464}\NoSourceHead
\pdfbookmark[1]{Zur sphärischen Trigonometrie}{spherical-trigonometry}
\HistoricTitle{[ZUR SPHÄRISCHEN TRIGONOMETRIE.]}
\indent \{Die Fundamentalgleichung
\[
\cos c\cos A=\sin c\cot b-\sin A\cot B
\]
\noindent kann, mit Vertauschung der Buchstaben, noch auf folgende 5 Arten geschrieben werden:
\[
\begin{aligned}
\cos a\cos C&=\sin a\cot b-\sin C\cot B\\
\cos c\cos B&=\sin c\cot a-\sin B\cot A\\
\cos b\cos C&=\sin b\cot a-\sin C\cot A\\
\cos b\cos A&=\sin b\cot c-\sin A\cot C\\
\cos a\cos B&=\sin a\cot c-\sin B\cot C.
\end{aligned}
\]

\indent Die Form, in welcher diese Gleichungen hier geschrieben sind, eignet sich am besten dazu, um sie im Gedächtniss zu behalten. Denn da die vier Bestandtheile, welche hier in Betracht kommen, im Umfange des Dreiecks unmittelbar auf einander folgen, so kann man immer zwei derselben wie mittlere, und die beiden andern wie äussere ansehen. Alsdann sagt die Gleichung in Worten: Das Product der Cosinus der \textls[120]{beiden mittleren} Bestandteile ist gleich dem Producte aus dem Sinus der \textls[120]{mittleren Seite} mit dem Cotangens der \textls[120]{äusseren Seite}, vermindert um das Product aus dem Sinus des \textls[120]{mittleren Winkels} mit dem Cotangens des \textls[120]{äusseren Winkels}. Man beachte dabei die Ähnlichkeit, welche das Vorkommen von cos, sin und cot mit den \textsc{Napier}schen Regeln hat.

\indent Dieses mnemonische Hülfsmittel pflegte \textsc{Gauss} seinen Zuhörern mitzuteilen.\}
\par\vfill\noindent X\textsubscript{1}.\hfill 58
\clearpage\SourceLeaf{465}\ArithmeticHead{458}{L}{GEOMETRIE. NACHLASS.}
\begingroup\small
\HistoricTitle{BEMERKUNG.}
\indent In dem \textit{Lehrbuch der Elementar-Mathematik} von \textsc{Th. Wittstein}, 2. Band, 2. Abteilung: \textit{Stereometrie}, Hannover 1862, finden sich, wie im Vorwort, S. VIII, hervorgehoben wird, drei auf \textsc{Gauss} zurückgehende Stellen. Die auf inhaltsgleiche sphärische Dreiecke bezügliche Stelle, S. 119, deckt sich mit dem Werke VIII, S. 293 abgedruckten Briefe von \textsc{Gauss} an \textsc{Schumacher} vom 6. Jan. 1842. Die in der zweiten, S. 161, gegebene Herleitung der sogenannten \textsc{Mollweide}schen Gleichungen findet sich schon in dem Werke VIII, S. 289 abgedruckten Briefe von \textsc{Gauss} an \textsc{Gerling} vom 18. Februar 1815. Dagegen verdient die S. 146 des \textsc{Wittstein}schen Buches mitgeteilte Gedächtnisregel für die Fundamentalgleichung hier nachgetragen zu werden.
\par\noindent\hfill\textsc{Stäckel}.
\ClosingRule
\par\endgroup
